% Pursuing Stacks (cotes 134-2 to 134-8): archive cote 134-6, site page 13 (PDF page 14), typescript page 332.
% typescript. Draft compared: ps5.tex:1989-2057.
% AI-transcribed by Claude Opus 5.5 (claude-opus-5-5, adaptive thinking, effort high), Message Batches API, batch ps3 (msgbatch_01XdFmuVVBAcgc7EFxdtCqsj), 2026-11-12; not yet checked by eye.
\scanpage{134-6/13}
\mspage{332}

it is immediately checked that Add(B) is an additive category \unsure{ahd|and} that
the functor (4) has the 2-universal property for functors from B into
any additive category, stated in yesterday's notes (p.330).

\uline{Remark}. The same construction essentially applies when considering the
universal problem of mapping B into any k-\uline{additive} category M \xout{\illegible{1}} (where
k is any \inserted[typed]{commutative} ring with unit), i.e. an additive category M endowed with a
ring homomorphism
\[ k \longrightarrow \mathrm{End}(\mathrm{id}_M) , \]
replacing \inserted[typed]{$\underline{Z}$ by k} in formulas (1) and (3) \xout{$\underline{Z}$xxbyxk}. The ``abelianization'' ques-%
tions touched at in yesterday's notes still make a sense in terms of
``k-linearization'' - a notion much in the spirit of our introduction
of a general basic localizer $\underline{W}$ , as the very notion of k-linearization
will give rise to a corresponding basic localizer $\underline{W}_k$ ...

Let's come back to the case when $B=A^o$ , A being a small category,
and to our question about chain complexes
\[ L. \quad \text{in} \quad Ch.(Add(A^o)) \]
giving rise to a commutative diagram (20) (p.328), up to isomorphism.
\over{a}{A} pair
\[ (L. , \hand{$\lambda$}) , \]
where L. is a chain complex as above, and \hand{$\lambda$} a compatibility isomor-%
phism for the square in (20), could be suggestively called an
\uline{abelianizator} for the small category A. The question of existence,
and uniqueness up to homotopy say, of an abelianizator for A seems
especially relevant when A is a test category say, and hence $A^\wedge$ mode-%
lizes homotopy types. In any case, in terms of an abelianizator we
get an additive functor
(5) \qquad $L\colon A^\wedge_{ab} \longrightarrow Ch.(Ab)$ ,
and the question arises whether this is compatible with weak equivalen-%
ces and quasi-isomorphisms\hand{;} \xout{-} maybe even,\xout{x} if \xout{it}\inserted{this} is not automatic, we
should insist it \xout{hs} holds when defining the notion of an abelianizator.
When this is OK, then by passing to localizations we deduce from (5)
a functor
(6) \qquad H\unsure{e|o}tab$_A \longrightarrow$ Hotab , \margin{\unsure{,}}
i.e. a functor (14) as looked for in yesterday's notes, giving rise
to \xout{a}\inserted[typed]{the} commutative diagram (19) (p.328) - whereas commutativity of
diagrams of the type \hand{(}15) (p. 326) looks less obvious.

\margin{13}
